\VerdierPhysicalPageAnchor{0068}%
\VerdierSemanticAnchor{ch1-2-2-1-continuation}%
groupes et que $\varepsilon(e_j,i)=-1$ si $j<i$,
$\varepsilon(e_j,i)=1$ si $i\leq j$. L'application
$(\sigma,i)\longmapsto\varepsilon(\sigma,i)$ est d'ailleurs uniquement
déterminée par ces propriétés. On en déduit immédiatement la relation :

\VerdierSemanticAnchor{ch1-2-2-1-eq2}
\begin{equation}
  \varepsilon(\sigma,i)\varepsilon(\sigma+e_i,j)
  =-\varepsilon(\sigma,j)\varepsilon(\sigma+e_j,i)
  \tag{2.2.1.2}
\end{equation}
valable pour tout $\sigma\in\mathbb Z[n]$ et pour tout couple
$(i,j)\in[n]\times[n]$, $i\neq j$.

\VerdierSemanticAnchor{ch1-2-2-2}
\noindent\textbf{2.2.2.}\quad
Soit $Y$ un complexe $n$-uple de $C$. Posons :

\VerdierSemanticAnchor{ch1-2-2-2-eq1}
\begin{equation}
  \delta_Y^{\sigma,i}=\varepsilon(\sigma,i)d_Y^{\sigma,i}.
  \tag{2.2.2.1}
\end{equation}

\VerdierSemanticAnchor{ch1-2-2-2-eq2}
On déduit alors de (2.2.1.2) que les relations (2.1.2.1) sont équivalentes
aux relations :
\begin{equation}
  \begin{aligned}
    \delta_Y^{\sigma+e_i,i}\delta_Y^{\sigma,i}
      &=0,
      &&\forall\sigma\in\mathbb Z[n],\quad\forall i\in[n],\\
    \delta_Y^{\sigma+e_j,i}\delta_Y^{\sigma,j}
      +\delta_Y^{\sigma+e_i,j}\delta_Y^{\sigma,i}
      &=0,
      &&\forall\sigma\in\mathbb Z[n],\quad\forall i,j\in[n].
  \end{aligned}
  \tag{2.2.2.2}
\end{equation}

\VerdierSemanticAnchor{ch1-2-2-2-pushforward}
Soit maintenant $\varphi\colon[n]\longrightarrow[m]$ une application et
supposons que $Y$ soit $\varphi$-sommable. On définit alors un complexe
$m$-uple $\int_\varphi Y$ par les formules :

\VerdierSemanticAnchor{ch1-2-2-2-eq3}
\begin{equation}
  \left(\int_\varphi Y\right)^\tau
  =\bigoplus_{\varphi(\sigma)=\tau}Y^\sigma,
  \qquad \forall\tau\in\mathbb Z[m],
  \tag{2.2.2.3}
\end{equation}

\VerdierSemanticAnchor{ch1-2-2-2-eq4}
\begin{equation}
  \begin{aligned}
    \delta_{\int_\varphi Y}^{\tau,i}(y_\sigma)
      &=\sum_{\varphi(j)=i}\delta_Y^{\sigma,j}(y_\sigma),\\[-1mm]
    &\hspace{-18mm}
      \sigma\in\mathbb Z[n],\quad i\in[m],\quad
      y_\sigma\in Y^\sigma,\quad\tau=\varphi(\sigma).
  \end{aligned}
  \tag{2.2.2.4}
\end{equation}

\VerdierSemanticAnchor{ch1-2-2-2-additive-interpretation}
La formule (2.2.2.4) définit un morphisme de
$\left(\int_\varphi Y\right)^\tau
=\bigoplus_{\varphi(\sigma)=\tau}Y^\sigma$ dans
$\left(\int_\varphi Y\right)^{\tau+e_i}
=\bigoplus_{\varphi(\sigma)=\tau+e_i}Y^\sigma$ lorsque $C$ est la
catégorie des groupes abéliens. Elle permet aussi de définir un morphisme
de $\left(\int_\varphi Y\right)^\tau$ dans
$\left(\int_\varphi Y\right)^{\tau+e_i}$ lorsque $C$ est une catégorie
additive en se plaçant dans la catégorie des foncteurs contravariants sur
$C$ à valeurs dans la catégorie des groupes abéliens, \emph{i.e.} en
interprétant $y_\sigma$ comme un morphisme d'un objet quelconque de $C$
dans $Y^\sigma$.

\VerdierSemanticAnchor{ch1-2-2-2-verification}
Un calcul immédiat montre que les
$\delta_{\int_\varphi Y}^{\tau,i}$ vérifient les relations (2.2.2.2)
d'où, en revenant aux différentielles, que les objets définis par (2.2.2.3)
munis des morphismes définis par :

\VerdierSemanticAnchor{ch1-2-2-2-eq5}
\begin{equation}
  \begin{aligned}
    d_{\int_\varphi Y}^{\tau,i}(y_\sigma)
      &=\varepsilon(\tau,i)
        \sum_{\varphi(j)=i}
          \varepsilon(\sigma,j)d_Y^{\sigma,j}(y_\sigma),\\[-1mm]
    &\hspace{-18mm}
      \tau\in\mathbb Z[m],\quad i\in[m],\quad
      \sigma\in\mathbb Z[n],\quad\varphi(\sigma)=\tau
  \end{aligned}
  \tag{2.2.2.5}
\end{equation}
